\ifdefined\gkver\else\def\gkver{student}\fi
\documentclass[\gkver]{gaokaozhenti}
\begin{document}

\chapter{2024年天津}

\begin{enumerate}
\item
%% number: 1
%% paperName: 2024年天津高考第 1 题 5 分
%% typeId: 1
%% type: 单选题
%% chapter: 光学
%% point: 全反射
%% method: 无
%% score: 5
%% degree: 950
%% duplicateId: 0
%% body:
下列关于光的说法正确的是\xzanswer{C}

\fourchoices[answer=C]
{光电效应揭示了光具有波动性}
{光从真空进入介质，波长不变}
{光纤通信利用了光的全反射现象}
{油膜呈现彩色条纹是光的偏振现象}

%% answer: C
%% memo:
\memoanswer{
A．光电效应揭示了光具有粒子性，故 A 错误；

B．光从真空进入介质，频率不变，波速减小，根据 $v = \lambda f$，可知波长变短，故 B 错误；

C．光纤通信利用了光的全反射现象，故 C 正确；

D．油膜呈现彩色条纹是光的薄膜干涉现象，故 D 错误。

故选 C。
}

\item
%% number: 2
%% paperName: 2024年天津高考第 2 题 5 分
%% typeId: 1
%% type: 单选题
%% chapter: 气体性质
%% point: 热力学第一定律
%% method: 无
%% score: 5
%% degree: 750
%% duplicateId: 0
%% body:
压缩空气储能是一种新型储能技术，其中涉及了空气的绝热膨胀和等温压缩过程，对一定质量的理想气体，下列说法正确的是\xzanswer{D}

\fourchoices[answer=D]
{绝热膨胀过程，气体的内能增大}
{绝热膨胀过程，外界对气体做功}
{等温压缩过程，气体的压强不变}
{等温压缩过程，气体的内能不变}

%% answer: D
%% memo:
\memoanswer{
AB．绝热膨胀过程，理想气体与外界没有热量的交换，即 $Q = 0$，理想气体对外界做功，即 $W < 0$，由热力学第一定律 $\Delta U = W + Q$，可知 $\Delta U < 0$，即气体的内能减小，故 AB 错误；

C．等温压缩过程，由玻意耳定律 $pV = C$，可知气体的压强变大，故 C 错误；

D．等温压缩过程，气体的温度保持不变，则内能不变，故 D 正确。

故选 D。
}

\item
%% number: 3
%% paperName: 2024年天津高考第 3 题 5 分
%% typeId: 1
%% type: 单选题
%% chapter: 机械波
%% point: 波与振动图象
%% method: 无
%% score: 5
%% degree: 650
%% duplicateId: 0
%% body:
一列简谐横波在均匀介质中沿 $x$ 轴传播，图 1 是 $t = 1 \Us$ 时该波的波形图，图 2 是 $x = 0$ 处质点的振动图像。则 $t = 11 \Us$ 时该波的波形图为\xzanswer{C}

\onepicture[w=9.00cm]{figs/03.png}

\fourchoices[answer=C, ispicture=true, h=2.60cm]
{figs/03a.png}
{figs/03b.png}
{figs/03c.png}
{figs/03d.png}

%% answer: C
%% memo:
\memoanswer{
波的周期 $T = 4 \Us$，因 $t = 11 \Us$ 时，即在 $t = 1 \Us$ 后再经过 $10 \Us = 2.5T$，此时原点处的质点振动到波谷位置，即该波的波形图为 C。

故选 C。
}

\item
%% number: 4
%% paperName: 2024年天津高考第 4 题 5 分
%% typeId: 1
%% type: 单选题
%% chapter: 电磁感应
%% point: 电磁感应中图象
%% method: 无
%% score: 5
%% degree: 550
%% duplicateId: 0
%% body:
如图所示，两根不计电阻的光滑金属导轨平行放置，导轨及其构成的平面均与水平面成某一角度，导轨上端用直导线连接，整个装置处在垂直于导轨平面向上的匀强磁场中。具有一定阻值的金属棒 MN 从某高度由静止开始下滑，下滑过程中 MN 始终与导轨垂直并接触良好，则 MN 所受的安培力 $F$ 及其加速度 $a$、速度 $v$、电流 $I$，随时间 $t$ 变化的关系图像可能正确的是\xzanswer{A}

\onepicture[w=6.00cm]{figs/04.png}

\fourchoices[answer=A, ispicture=true, h=2.00cm]
{figs/04a.png}
{figs/04b.png}
{figs/04c.png}
{figs/04d.png}

%% answer: A
%% memo:
\memoanswer{
ABC．根据题意，设导体棒的电阻为 $R$，导轨间距为 $L$，磁感应强度为 $B$，导体棒速度为 $v$ 时，受到的安培力为

$F = BIL = \frac{B^{2}L^{2}v}{R}$

可知 $F \propto v$

由牛顿第二定律可得，导体棒的加速度为

$a = \frac{mg\sin\theta - F}{m} = g\sin\theta - \frac{B^{2}L^{2}v}{mR}$

可知，随着速度的增大，导体棒的加速度逐渐减小，当加速度为零时，导体棒开始做匀速直线运动，则 $v-t$ 图像的斜率逐渐减小直至为零时，速度保持不变；由于安培力 $F$ 与速度 $v$ 成正比，则 $F-t$ 图像的斜率逐渐减小直至为零时，$F$ 保持不变，故 A 正确，BC 错误；

D．根据题意，由公式可得，感应电流为 $I = \frac{BLv}{R}$

由数学知识可得 $\frac{\Delta I}{\Delta t} = \frac{BL}{R}\cdot\frac{\Delta v}{\Delta t} = \frac{BLa}{R}$

由于加速度逐渐减小，则 $I-t$ 图像的斜率逐渐减小，故 D 错误。

故选 A。
}

\item
%% number: 5
%% paperName: 2024年天津高考第 5 题 5 分
%% typeId: 1
%% type: 单选题
%% chapter: 运动和力
%% point: 牛顿第二定律
%% method: 无
%% score: 5
%% degree: 550
%% duplicateId: 0
%% body:
生活中人们经常使用如图所示的小车搬运重物，小车的底板和侧板垂直，底板和侧板对重物的弹力分别为 $F_{\text{底}}$、$F_{\text{侧}}$，忽略重物和底板之间的摩擦力。保持底板与水平面之间的夹角不变，重物始终与小车相对静止，小车水平向右运动，由匀速变为加速时\xzanswer{B}

\onepicture[w=4.00cm]{figs/05.png}

\fourchoices[answer=B]
{$F_{\text{底}}$ 增大，$F_{\text{侧}}$ 增大}
{$F_{\text{底}}$ 减小，$F_{\text{侧}}$ 增大}
{$F_{\text{底}}$ 增大，$F_{\text{侧}}$ 减小}
{$F_{\text{底}}$ 减小，$F_{\text{侧}}$ 减小}

%% answer: B
%% memo:
\memoanswer{
对重物受力分析，设底板与水平面之间的夹角为 $\theta$，如图所示

小车匀速时，有 $F_{\text{底}}\sin\theta = F_{\text{侧}}\cos\theta$；小车加速时，有

$F_{\text{侧}}\cos\theta - F_{\text{底}}\sin\theta = ma$

其中 $a > 0$，则 $F_{\text{底}}$ 减小，$F_{\text{侧}}$ 增大。

故选 B。

\onepicture[w=3.00cm]{figs/05an.png}
}

\item
%% number: 6
%% paperName: 2024年天津高考第 6 题 5 分
%% typeId: 2
%% type: 多选题
%% chapter: 原子和原子核
%% point: 人工核反应
%% method: 无
%% score: 5
%% degree: 750
%% duplicateId: 0
%% body:
中国钍基熔盐堆即将建成小型实验堆，为我国能源安全和可持续发展提供有力支持。反应堆中涉及的核反应方程有：①$\ce{X + ^{232}_{90}Th -> ^{233}_{90}Th}$，②$\ce{^{233}_{90}Th -> ^{233}_{91}Pa + ^{0}_{-1}e}$，下列说法正确的是\xzanswer{AB}

\fourchoices[answer=AB]
{方程①中 X 是中子}
{方程②中 $\ce{^{233}_{90}Th}$ 发生了 $\beta$ 衰变}
{受反应堆高温影响，$\ce{^{233}_{90}Th}$ 的半衰期会变短}
{方程②释放电子，说明电子是原子核的组成部分}

%% answer: AB
%% memo:
\memoanswer{
A．根据核反应的质量数和电荷数守恒可知，方程①中 X 质量数为 1，电荷数为零，则是中子，选项 A 正确；

B．方程②中 $\ce{^{233}_{90}Th}$ 放出负电子，则发生了 $\beta$ 衰变，选项 B 正确；

C．放射性元素的半衰期与外界因素无关，选项 C 错误；

D．方程②释放的电子是原子核内的中子转化为质子时放出的，不能说明电子是原子核的组成部分，选项 D 错误。

故选 AB。
}

\item
%% number: 7
%% paperName: 2024年天津高考第 7 题 5 分
%% typeId: 2
%% type: 多选题
%% chapter: 曲线运动
%% point: 人造地球卫星
%% method: 无
%% score: 5
%% degree: 550
%% duplicateId: 0
%% body:
卫星未发射时静置在赤道上随地球转动，地球半径为 $R$。卫星发射后在地球同步轨道上做匀速圆周运动，轨道半径为 $r$。则卫星未发射时和在轨道上运行时\xzanswer{AC}

\fourchoices[answer=AC]
{角速度之比为 $1:1$}
{线速度之比为 $\sqrt{r}:\sqrt{R}$}
{向心加速度之比为 $R:r$}
{受到地球的万有引力之比为 $r^{2}:R^{2}$}

%% answer: AC
%% memo:
\memoanswer{
A．卫星未发射时静置在赤道上随地球转动，角速度与地球自转角速度相等，卫星发射后在地球同步轨道上做匀速圆周运动，角速度与地球自转角速度相等，则卫星未发射时和在轨道上运行时角速度之比为 $1:1$，故 A 正确；

B．根据题意，由公式 $v = \omega r$ 可知，卫星未发射时和在轨道上运行时，由于角速度相等，则线速度之比为轨道半径之比 $R:r$，故 B 错误；

C．根据题意，由公式 $a_{n} = \omega^{2}r$ 可知，卫星未发射时和在轨道上运行时，由于角速度相等，则向心加速度之比为轨道半径之比 $R:r$，故 C 正确；

D．根据题意，由万有引力定律可知，卫星未发射时和在轨道上运行时，受到地球的万有引力之比与轨道半径的平方成反比，即 $r^{2}:R^{2}$，故 D 错误。

故选 AC。
}

\item
%% number: 8
%% paperName: 2024年天津高考第 8 题 5 分
%% typeId: 2
%% type: 多选题
%% chapter: 电场
%% point: 电场中的图象
%% method: 无
%% score: 5
%% degree: 450
%% duplicateId: 0
%% body:
某静电场在 $x$ 轴正半轴的电势 $\varphi$ 随 $x$ 变化的图像如图所示，a、b、c、d 为 $x$ 轴上四个点。一负电荷仅在静电力作用下，以一定初速度从 d 点开始沿 $x$ 轴负方向运动到 a 点，则该电荷\xzanswer{BD}

\onepicture[w=4.50cm]{figs/08.png}

\fourchoices[answer=BD]
{在 b 点电势能最小}
{在 c 点时速度最小}
{所受静电力始终做负功}
{在 a 点受静电力沿 $x$ 轴负方向}

%% answer: BD
%% memo:
\memoanswer{
AB．根据题意，由公式 $E_{p} = q\varphi$ 可知，负电荷在高电势位置的电势能较小，由图可知，a 点的电势最大，则在 a 点电势能最小，同理可知，c 点的电势最小，则在 c 点时电势能最大，电荷仅在电场力作用下，电荷的电势能和动能之和不变，可知，电势能最大时，动能最小，则在 c 点时，电荷的动能最小，即速度最小，故 A 错误，B 正确；

CD．根据沿电场线方向电势逐渐降低，结合题图可知，c 点左侧电场方向沿 $x$ 轴正方向，c 点右侧电场方向沿 $x$ 轴负方向，可知，c 点右侧负电荷受沿 $x$ 轴正方向的电场力，c 点左侧负电荷受沿 $x$ 轴负方向的电场力，可知，在 a 点受静电力沿 $x$ 轴负方向，从 d 点开始沿 $x$ 轴负方向运动到 a 点，电场力先做负功再做正功，故 C 错误，D 正确。

故选 BD。
}

\item
%% number: 9
%% paperName: 2024年天津高考第 9 题 6 分
%% typeId: 4
%% type: 实验题
%% chapter: 电路
%% point: 闭合电路欧姆定律
%% method: 无
%% score: 6
%% degree: 650
%% duplicateId: 0
%% body:
某同学研究闭合电路的规律。
\begin{enumerate}
	\item
	根据闭合电路的欧姆定律得出了电源输出功率 $P$ 与外电路电阻关系图像，如图所示，则 $P$ 的峰值对应的外电路电阻值 $R$ 应\tkanswer{等于}电源内阻 $r$（填“大于”、“小于”或“等于”）；
	\item
	测定电源的电动势和内阻，可供选用的器材有：

	A．电压表：（量程 $0 \sim 3 \UV$，内阻约为 $3 \UkO$）

	B．电流表：（量程 $0 \sim 0.6 \UA$，内阻约为 $1 \UO$）

	C．滑动变阻器：（最大阻值 $20 \UO$，额定电流 $1 \UA$）

	D．滑动变阻器：（最大阻值 $1000 \UO$，额定电流 $0.5 \UA$）

	E．待测电源：（电动势约为 $3 \UV$，内阻约为 $1 \UO$）

	F．开关、导线若干

	（i）实验中所用的滑动变阻器应选\tkanswer{C}（填器材前字母代号）；

	（ii）实物电路如图所示，单刀双掷开关 S$_{2}$ 可分别与 1、2 端闭合，为使电源内阻的测量结果更接近真实值，S$_{2}$ 应与\tkanswer{2}端闭合。
\end{enumerate}
\twopicture[wA=3.20cm, wB=4.60cm, align=right]
{figs/09a.png}
{figs/09b.png}
%% answer:
\jdanswer{
\begin{enumerate}
	\item
	等于

	\item
	C，2

\end{enumerate}
}
%% memo:
\memoanswer{
（1）电源输出功率

$P = \frac{E^{2}R}{(R + r)^{2}} = \frac{E^{2}}{\frac{(R - r)^{2}}{R} - 4r}$

则当 $R = r$ 时电源输出功率 $P$ 最大；

（2）（i）实验中所用的滑动变阻器应选阻值较小的 C 即可；

（ii）电压表内阻远大于电源内阻，应采用相对电源的电流表外接法，使电源内阻的测量结果更接近真实值，S$_{2}$ 应与 2 端闭合。
}

\item
%% number: 9
%% paperName: 2024年天津高考第 9 题 6 分
%% typeId: 4
%% type: 实验题
%% chapter: 运动和力
%% point: 验证牛顿第二定律
%% method: 无
%% score: 6
%% degree: 550
%% duplicateId: 0
%% body:
某同学用图示装置探究加速度与力的关系。
\onepicture[w=6.50cm, align=right]{figs/09c.png}
\begin{enumerate}
	\item
	为补偿打点计时器对小车的阻力及其他阻力，调节木板倾角，使小车在不挂槽码时运动，并打出纸带进行检验，下图中能表明补偿阻力恰当的是\tkanswer{B}；

	\fourchoices[answer=B, ispicture=true, h=0.80cm]
	{figs/09d.png}
	{figs/09e.png}
	{figs/09f.png}
	{figs/09g.png}
	\item
	某次实验得到一条纸带，部分计数点如下图所示（每相邻两个计数点间还有 4 个点，图中未画出），测得 $s_{1} = 6.20 \Ucm$，$s_{2} = 6.70 \Ucm$，$s_{3} = 7.21 \Ucm$，$s_{4} = 7.73 \Ucm$。已知打点计时器所接交流电源频率为 $50 \UHz$，则小车的加速度 $a = $\tkanswer{0.51}$\Umsq$（要求充分利用测量数据，结果保留两位有效数字）；
	\onepicture[w=8.00cm, align=right]{figs/09h.png}
	\item
	该同学将一个可以直接测出绳子拉力的传感器安装在小车上，小车和传感器总质量为 $210 \Ug$。按要求补偿阻力后，该同学共进行了四次实验，悬挂的槽码质量依次为 $5 \Ug$、$10 \Ug$、$20 \Ug$、$40 \Ug$。处理数据时，用两种方式得到小车（含传感器）受到的合力，一种将槽码所受重力当作合力、另一种将传感器示数当作合力，则这两种方式得到的合力差异最大时，槽码质量为\tkanswer{40} g。
\end{enumerate}
%% answer:
\jdanswer{
\begin{enumerate}
	\item
	B

	\item
	0.51

	\item
	40

\end{enumerate}
}
%% memo:
\memoanswer{
（1）若补偿摩擦力恰当，则小车应该匀速运动，打出的纸带应该点迹均匀分布，故选 B。

（2）每相邻两个计数点间还有 4 个点未画出，可知 $T = 0.1 \Us$；小车的加速度

$a = \frac{s_{4} + s_{3} - s_{2} - s_{1}}{4T^{2}} = \frac{(7.73 + 7.21 - 6.70 - 6.20)\times10^{-2}}{4\times0.1^{2}}\Umsq = 0.51\Umsq$

（3）根据牛顿第二定律，对砝码 $mg - T = ma$

对小车 $T = Ma$

可得 $T = \frac{Mmg}{M + m} = \frac{1}{1 + \frac{m}{M}}mg$

则当 $m$ 较小时传感器的示数越接近砝码的重力 $mg$；$m$ 越大，则传感器的示数与砝码重力的差异越大，即这两种方式得到的合力差异最大时，槽码质量为 $40 \Ug$。
}

\item
%% number: 10
%% paperName: 2024年天津高考第 10 题 14 分
%% typeId: 6
%% type: 计算题
%% chapter: 运动和力
%% point: 动量守恒定律
%% method: 无
%% score: 14
%% degree: 550
%% duplicateId: 0
%% body:
如图所示，光滑半圆轨道直径沿竖直方向，最低点与水平面相切。对静置于轨道最低点的小球 A 施加水平向左的瞬时冲量 $I$，A 沿轨道运动到最高点时，与用轻绳悬挂的静止小球 B 正碰并粘在一起。已知 $I = 1.8 \UNs$，A、B 的质量分别为 $m_{A} = 0.3 \Ukg$、$m_{B} = 0.1 \Ukg$，轨道半径和绳长均为 $R = 0.5 \Um$，两球均视为质点，轻绳不可伸长，重力加速度 $g$ 取 $10 \Umsq$，不计空气阻力。求：
\begin{enumerate}
	\item
	与 B 碰前瞬间 A 的速度大小；
	\item
	A、B 碰后瞬间轻绳的拉力大小。
\end{enumerate}
\onepicture[w=3.50cm, align=right]{figs/10.png}
%% answer:
\jdanswer{
\begin{enumerate}
	\item
	$4 \Ums$

	\item
	$11.2 \UN$

\end{enumerate}
}
%% memo:
\memoanswer{
（1）根据题意，设小球 A 从最低点开始运动时的速度为 $v_{0}$，由动量定理有 $I = m_{A}v_{0}$

设与 B 碰前瞬间 A 的速度大小 $v$，从最低点到最高点，由动能定理有

$-mg\cdot 2R = \frac{1}{2}mv^{2} - \frac{1}{2}mv_{0}^{2}$

联立代入数据解得 $v = 4 \Ums$

（2）A 与用轻绳悬挂的静止小球 B 正碰并粘在一起，由动量守恒定律有

$m_{A}v = (m_{A} + m_{B})v_{\text{共}}$

设 A、B 碰后瞬间轻绳的拉力大小为 $F$，由牛顿第二定律有

$F - (m_{A} + m_{B})g = (m_{A} + m_{B})\frac{v_{\text{共}}^{2}}{R}$

联立代入数据解得 $F = 11.2 \UN$
}

\item
%% number: 11
%% paperName: 2024年天津高考第 11 题 16 分
%% typeId: 6
%% type: 计算题
%% chapter: 磁场
%% point: 带电粒子在磁场中的运动
%% method: 无
%% score: 16
%% degree: 450
%% duplicateId: 0
%% body:
如图所示，在 $Oxy$ 平面直角坐标系的第一象限内，存在半径为 $R$ 的半圆形匀强磁场区域，半圆与 $x$ 轴相切于 M 点，与 $y$ 轴相切于 N 点，直线边界与 $x$ 轴平行，磁场方向垂直于纸面向里。在第一象限存在沿 $+x$ 方向的匀强电场，电场强度大小为 $E$。一带负电粒子质量为 $m$，电荷量为 $q$，从 M 点以速度 $v$ 沿 $+y$ 方向进入第一象限，正好能沿直线匀速穿过半圆区域。不计粒子重力。
\begin{enumerate}
	\item
	求磁感应强度 $B$ 的大小；
	\item
	若仅有电场，求粒子从 M 点到达 $y$ 轴的时间 $t$；
	\item
	若仅有磁场，改变粒子入射速度的大小，粒子能够到达 $x$ 轴上 P 点，M、P 的距离为 $\sqrt{3}R$，求粒子在磁场中运动的时间 $t_{1}$。
\end{enumerate}
\onepicture[w=5.00cm, align=right]{figs/11.png}
%% answer:
\jdanswer{
\begin{enumerate}
	\item
	$B = \frac{E}{v}$

	\item
	$t = \sqrt{\frac{2mR}{qE}}$

	\item
	$t_{1} = \frac{2\pi mv}{3qE}$

\end{enumerate}
}
%% memo:
\memoanswer{
（1）根据题意可知，由于一带负电粒子能沿直线匀速穿过半圆区域，由平衡条件有 $qE = qvB$

解得 $B = \frac{E}{v}$

（2）若仅有电场，带负电粒子受沿 $x$ 轴负方向的电场力，由牛顿第二定律有 $qE = ma$

又有 $R = \frac{1}{2}at^{2}$

联立解得 $t = \sqrt{\frac{2mR}{qE}}$

（3）根据题意，设粒子入射速度为 $v_{0}$，则有

$qv_{0}B = m\frac{v_{0}^{2}}{r}$，$T = \frac{2\pi r}{v_{0}}$

可得 $T = \frac{2\pi m}{qB} = \frac{2\pi mv}{qE}$

画出粒子的运动轨迹，如图所示

由几何关系可得 $\tan\theta = \frac{\sqrt{3}R}{R} = \sqrt{3}$

解得 $\theta = 60^\circ$

则轨迹所对圆心角为 $120^\circ$，则粒子在磁场中运动的时间

$t_{1} = \frac{120^\circ}{360^\circ}T = \frac{2\pi mv}{3qE}$

\onepicture[w=4.50cm]{figs/11an.png}
}

\item
%% number: 12
%% paperName: 2024年天津高考第 12 题 18 分
%% typeId: 6
%% type: 计算题
%% chapter: 电磁感应
%% point: 电磁感应能量分析
%% method: 无
%% score: 18
%% degree: 350
%% duplicateId: 0
%% body:
电动汽车制动过程中可以控制电机转为发电模式，在产生制动效果的同时，将汽车的部分机械能转换为电能，储存在储能装置中，实现能量回收、降低能耗。如图 1 所示，发电机可简化为处于匀强磁场中的单匝正方形线框 ABCD，线框边长为 $L$，电阻忽略不计，磁场磁感应强度大小为 $B$，线框转轴 OO$'$ 与磁场垂直，且与 AB、CD 距离相等。线框与储能装置连接。
\begin{enumerate}
	\item
	线框转动方向如图 1 所示，试判断图示位置 AB 中的电流方向；
	\item
	若线框以角速度 $\omega$ 匀速转动，线框平面与中性面垂直瞬间开始计时，线框在 $t$ 时刻位置如图 2 所示，求此时 AB 产生的感应电动势；
	\item
	讨论电动汽车在某次制动储存电能时，为方便计算，做两点假定：①将储能装置替换为阻值为 $R$ 的电阻，电阻消耗的电能等于储能装置储存的电能；②线框转动第一周的角速度为 $\omega_{0}$，第二周的角速度为 $\frac{\omega_{0}}{2}$，第三周的角速度为 $\frac{\omega_{0}}{4}$，依次减半，直到线框停止转动。若该制动过程中汽车在水平路面上做匀减速直线运动，汽车质量为 $m$，加速度大小为 $a$，储存的电能为初动能的 50\%，求制动过程中汽车行驶的最大距离 $x$。
\end{enumerate}
\onepicture[w=9.00cm, align=right]{figs/12.png}
%% answer:
\jdanswer{
\begin{enumerate}
	\item
	电流方向从 B 到 A

	\item
	$e = \frac{1}{2}\omega BL^{2}\cos\omega t$

	\item
	$x = \frac{4\pi\omega_{0}B^{2}L^{4}}{maR}$

\end{enumerate}
}
%% memo:
\memoanswer{
（1）由右手定则可知，图示位置 AB 中的电流方向从 B 到 A。

（2）线框平面与中性面垂直瞬间开始计时，则经过时间 $t$ 转过的角度为 $\theta = \omega t$

AB 切割磁感线的速度为 $v = \omega\frac{L}{2}\cos\theta$

则感应电动势 $e = BLv$

解得此时 AB 产生的感应电动势 $e = \frac{1}{2}\omega BL^{2}\cos\omega t$

（3）线框转动过程中，AB、CD 均能产生感应电动势，故线框转动产生的感应电动势为 $e_{1} = 2e = \omega BL^{2}\cos\omega t$

线框转动第一周产生感应电动势最大值 $E_{m1} = B\omega_{0}L^{2}$

储存电能为 $\varepsilon_{1} = \frac{\left(\frac{E_{m1}}{\sqrt{2}}\right)^{2}}{R}\cdot\frac{2\pi}{\omega_{0}} = \frac{\pi B^{2}L^{4}\omega_{0}}{R}$

同理线框转动第二周储存的电能 $\varepsilon_{2} = \frac{\pi B^{2}L^{4}}{R}\cdot\frac{\omega_{0}}{2}$

同理线框转动第三周储存的电能 $\varepsilon_{3} = \frac{\pi B^{2}L^{4}}{R}\cdot\frac{\omega_{0}}{2^{2}}$

……

线框转动第 $n$ 周储存的电能 $\varepsilon_{n} = \frac{\pi B^{2}L^{4}}{R}\cdot\frac{\omega_{0}}{2^{n-1}}$

则直到停止时储存的电能为

$\varepsilon = \varepsilon_{1} + \varepsilon_{2} + \cdots + \varepsilon_{n} = \frac{\pi B^{2}L^{4}\omega_{0}}{R}\left(1 + \frac{1}{2} + \cdots + \frac{1}{2^{n-1}}\right) = \frac{2\pi B^{2}L^{4}\omega_{0}}{R}$

储存的电能为初动能的 50\%，可知初动能 $E_{k0} = \frac{4\pi B^{2}L^{4}\omega_{0}}{R}$

根据动能定理和牛顿第二定律可得 $fx = E_{k0} - 0$，$f = ma$

解得 $x = \frac{4\pi\omega_{0}B^{2}L^{4}}{maR}$
}

\end{enumerate}

\end{document}
